\documentclass[11pt]{article}
\usepackage[T1]{fontenc}
\usepackage{lmodern,microtype,amsmath,amssymb,amsthm,booktabs}
\microtypesetup{expansion=false}
\usepackage[margin=1.05in,headheight=14pt]{geometry}
\usepackage{xcolor}
\definecolor{linknavy}{RGB}{25,55,125}
\usepackage[colorlinks=true,linkcolor=linknavy,citecolor=linknavy,urlcolor=linknavy]{hyperref}
\usepackage{orcidlink,fancyhdr}
\hypersetup{
  pdftitle={How much dissipation can a multiscale layer cascade carry},
  pdfauthor={Felipe Santibanez-Leal},
  pdfsubject={Upper bounds on the fractional dissipation exponent compatible with layer-cascade blowup constructions},
  pdfkeywords={Navier-Stokes, Boussinesq, finite-time blowup, hypodissipative, multiscale construction, fractional Laplacian}
}
\newtheorem{theorem}{Theorem}[section]
\newtheorem{proposition}[theorem]{Proposition}
\newtheorem{corollary}[theorem]{Corollary}
\newtheorem{lemma}[theorem]{Lemma}
\theoremstyle{remark}
\newtheorem{remark}[theorem]{Remark}
\newcommand{\Q}{\mathbb Q}
\newcommand{\versiondoi}{10.5281/zenodo.22821790}
\newcommand{\conceptdoi}{10.5281/zenodo.22820520}
\newcommand{\recordurl}{https://github.com/fsantibanezleal/CAOS_RESEARCH/tree/main/problems/analysis-pde/navier-stokes}
\pagestyle{fancy}
\fancyhf{}
\fancyhead[L]{\small CAOS Research Preprint}
\fancyhead[R]{\small Dissipation limits of layer cascades\textperiodcentered\ v0.04}
\fancyfoot[C]{\thepage}
\fancypagestyle{plain}{\fancyhf{}\fancyfoot[C]{\thepage}\renewcommand{\headrulewidth}{0pt}}
\allowdisplaybreaks[1]

\title{How much dissipation can a multiscale layer cascade carry?}
\author{Felipe Santiba\~nez-Leal\,\orcidlink{0000-0002-0150-3246}\\
\small CAOS Research program\\
\small\texttt{github.com/fsantibanezleal/CAOS\_RESEARCH}}
\date{2026-09-17\\[2pt]\normalsize Version 0.04}

\begin{document}
\maketitle
\thispagestyle{plain}
\begin{center}\small
\fbox{\parbox{0.92\textwidth}{\centering
\textbf{PREPRINT}\textperiodcentered\ not peer reviewed\textperiodcentered\ CAOS Research\\
Version 0.04\textperiodcentered\ 2026-09-17\textperiodcentered\ CC BY 4.0\\
Version DOI: \href{https://doi.org/\versiondoi}{\versiondoi}\\
Concept DOI: \href{https://doi.org/\conceptdoi}{\conceptdoi} (latest version)\\
\href{\recordurl}{Sources, numerical data, and machine-checked computations}
}}
\end{center}

\begin{abstract}
Multiscale layer constructions produce finite-time blowup for several fluid equations: for the
incompressible porous media equation, for the two-dimensional Boussinesq system, for
three-dimensional Euler, and, with fractional dissipation $|\nabla|^{\alpha}$ of order
$\alpha<(22-8\sqrt7)/9$, for the forced Navier-Stokes equations. The exponent in the last of these is
small, and it is natural to ask what limits it.

We give two upper bounds on the dissipation exponent such constructions can carry. The first is
independent of every design constant: a cascade whose layers grow at rate $\sqrt A$ on a background
gradient $A$ cannot exceed $\alpha=1/2$, and one whose layers grow at rate $A$ cannot exceed
$\alpha=1$, both in the $|\nabla|^{\alpha}$ convention in which classical viscosity is $\alpha=2$.
Neither mechanism reaches classical viscosity, by factors of four and two respectively. The second
bound applies to the smooth-forcing design used for the inviscid Boussinesq and Euler results: its
correction hierarchy, seed rule and frequency-ratio rule cap the exponent at
$\alpha\le1.08\times10^{-3}$, attained when exactly one derivative of the force is controlled, which
is $86$ times below the exponent already achieved with a force of finite regularity.

Both bounds follow from the amplitude system of the construction together with one additional
ingredient, the effect of fractional dissipation on a single layer, which we derive and then verify
against direct simulation of the nonlinear equations, including the effect of the spatial
localization the constructions use. We also reconstruct the exponent budget behind the published
hypodissipative threshold, identify the constraint that determines it, and show that the admissible
frequency ratios form an interval that closes exactly at that threshold, so that cascades with
geometrically growing frequencies are excluded at every positive dissipation.
\end{abstract}

\section{The question}

Let $\theta$ and $u$ solve the forced two-dimensional Boussinesq system, or let $u$ solve the forced
three-dimensional Navier-Stokes equations with fractional dissipation. Multiscale constructions in the
style introduced by C\'ordoba and Mart\'inez-Zoroa build blowup by stacking oscillatory layers: each
layer grows on the gradient deposited by its predecessors and hands a larger gradient to its
successor, with infinitely many stages compressed into finite time.

Dissipation opposes this directly. A layer of frequency $\lambda$ is damped at rate
$\nu\lambda^{2\alpha}$ in the convention $-\nu(-\Delta)^{\alpha}$, while its growth rate is set by the
background gradient alone. Frequency therefore buys gradient but not growth, and the question of how
much dissipation a cascade tolerates becomes a question about the exponents relating amplitude,
frequency and gradient. This paper answers it at two levels: for the mechanism classes
(Section~\ref{sec:class}) and for one published design (Section~\ref{sec:design}).

\paragraph{Conventions.} We write dissipation as $-\nu(-\Delta)^{\alpha}$, so classical viscosity is
$\alpha=1$; results quoted from the literature in the $|\nabla|^{\alpha}$ convention have twice our
exponent, and we state which convention is in use whenever a number appears. The published
hypodissipative threshold $(22-8\sqrt7)/9=0.09267$ in the $|\nabla|^{\alpha}$ convention is
$0.04633$ in ours.

\section{The layer mechanism and its dissipative extension}\label{sec:mechanism}

On a background that is affine near the origin, $u_{\mathrm{old}}=D(t)x$ and
$\theta_{\mathrm{old}}=G(t)\cdot x$, a temperature wave and a vorticity wave sharing one phase
$s=\lambda\zeta(t)\cdot x$, with $\vartheta=\Theta F(s)$, $\varpi=\Omega F'(s)$ and
$\psi=\Omega P(s)/(\lambda^{2}|\zeta|^{2})$, satisfy a closed system of amplitude equations,
\begin{equation}\label{eq:mod}
\dot\zeta=-D^{T}\zeta,\qquad
\dot\Theta=-\frac{J\zeta\cdot G}{\lambda|\zeta|^{2}}\,\Omega,\qquad
\dot\Omega=\lambda\,\zeta_{1}\,\Theta ,
\end{equation}
where $J(x_{1},x_{2})=(-x_{2},x_{1})$. The wave velocity is perpendicular to $\zeta$ while both wave
gradients are parallel to it, so the wave does not advect itself and the ansatz is exact where the
envelope is one and the profile is linear. In the frozen configuration $D=0$, $G=-Ae_{2}$,
$\zeta=re(\varphi)$, the growth rate is $\sqrt A\sin\varphi$: it carries no $\lambda$, because the
frequency in $\dot\Omega$ cancels the reciprocal frequency in $\dot\Theta$. The gradient the layer
deposits is governed by the amplification identity
\begin{equation}\label{eq:amp}
\nabla\vartheta(0,t)=\lambda\Theta\zeta ,
\end{equation}
which is why a short wave can carry a small amplitude and a large gradient, and is what makes the
cascade possible.

\paragraph{Dissipation.} Adding $-\nu(-\Delta)^{\alpha}$ to both equations leaves the two geometric
cancellations untouched and is diagonal on the phase, since
$(-\Delta)^{\alpha}F(\lambda\zeta\cdot x)=(\lambda|\zeta|)^{2\alpha}F(\lambda\zeta\cdot x)$. The
amplitude system therefore acquires one damping term per amplitude,
\begin{equation}\label{eq:diss}
\dot\Theta=-\frac{J\zeta\cdot G}{\lambda|\zeta|^{2}}\,\Omega-\nu(\lambda|\zeta|)^{2\alpha}\Theta,
\qquad
\dot\Omega=\lambda\zeta_{1}\Theta-\nu(\lambda|\zeta|)^{2\alpha}\Omega ,
\end{equation}
with eigenvalues $-\nu(\lambda r)^{2\alpha}\pm\sqrt A\sin\varphi$ in the frozen configuration. Growth
at frequency $\lambda$ on background gradient $A$ therefore requires
\begin{equation}\label{eq:growth}
\nu(\lambda r)^{2\alpha}<\sqrt A\,\sin\varphi .
\end{equation}
Equation~\eqref{eq:diss} is the one ingredient of this paper that does not appear in the
constructions themselves; Section~\ref{sec:validation} reports its verification against direct
simulation of the nonlinear equations.

\paragraph{Two growth laws.} The Boussinesq layers above grow at $\sqrt A$, a pendulum law: the
amplitude pair obeys a second-order equation whose coefficient is the background gradient. The
vortex layers of the hypodissipative Navier-Stokes construction instead grow at rate $A_{n-1}$, the
background gradient itself, which is why its dissipation constraint reads $a=\alpha R$ in the
exponent notation of Section~\ref{sec:budget} rather than carrying a square root. Both laws are
treated below.

\section{A bound independent of design constants}\label{sec:class}

\begin{theorem}\label{thm:class}
Consider a layer cascade with the following structure:
\begin{enumerate}
\item[(A1)] each layer stops at amplitude $\lambda_{q}^{-1+\delta}$ with $\delta<1$;
\item[(A2)] it deposits a background gradient $A_{q+1}=\lambda_{q}^{\delta}$, by~\eqref{eq:amp};
\item[(A3)] frequencies and background gradients are related by $\lambda_{q}=A_{q}^{p}$ for a fixed $p$;
\item[(A4)] layers grow at rate $A^{\gamma}$ with $\gamma=1/2$ or $\gamma=1$;
\item[(A5)] a layer of frequency $\lambda$ is damped at $\nu\lambda^{2\alpha}$, as in~\eqref{eq:diss}.
\end{enumerate}
Then $p\delta>1$, and growth at every stage requires
\[
\alpha\le\frac{\gamma}{2p}<\frac{\gamma\,\delta}{2}<\frac{\gamma}{2},
\]
that is $\alpha<1/4$ for $\gamma=1/2$ and $\alpha<1/2$ for $\gamma=1$. In the $|\nabla|^{\alpha}$
convention these ceilings are $1/2$ and $1$, against a classical viscosity of $2$.
\end{theorem}

\begin{proof}
By (A2) and (A3), $A_{q+1}=\lambda_{q}^{\delta}=A_{q}^{p\delta}$. The gradients must grow without bound
for the cascade to produce blowup, so $p\delta>1$; with $\delta<1$ from (A1), which is required because
the field itself stays bounded while its gradient diverges, this gives $p>1/\delta>1$.

By (A5) and (A4), stage $q$ grows only while $\nu\lambda_{q}^{2\alpha}<c\,A_{q}^{\gamma}$ for a
constant $c$ independent of $q$. Taking logarithms and using (A3),
\[
\log\nu+2\alpha p\log A_{q}<\log c+\gamma\log A_{q},
\]
and since $A_{q}\to\infty$ the inequality can hold for all large $q$ only if $2\alpha p\le\gamma$,
that is $\alpha\le\gamma/(2p)$. At equality the condition reduces to $\nu<c$, which a small viscosity
satisfies, so this step is not strict; strictness of the final bounds comes from $p>1/\delta$, which
gives $\alpha\le\gamma/(2p)<\gamma\delta/2<\gamma/2$.
\end{proof}

In the variables of Section~\ref{sec:design}, where frequencies satisfy
$\lambda_{q}=\lambda_{q-1}^{Q}$, the exponents are related by $Q=p\delta$, and the bound
$\alpha\le\gamma/(2p)$ reads $\alpha\le\gamma\delta/(2Q)$, which for $\gamma=1/2$ is the constraint
$\alpha\le\delta/(4Q)$ used there.

\begin{corollary}\label{cor:classical}
No cascade of either type reaches classical viscosity. The pendulum mechanism falls short by a factor
of four and the vortex-stretching mechanism by a factor of two, uniformly in every constant of the
design.
\end{corollary}

The two ceilings differ by exactly the growth law: a layer that grows at the background gradient,
rather than at its square root, gains a factor of two and nothing else in the argument changes. The
published constructions sit well below their own class ceilings, the hypodissipative vortex-layer
threshold by a factor $10.8$ and the design of Section~\ref{sec:design} by a factor $465$.

\begin{remark}
Corollary~\ref{cor:classical} is consistent with a choice visible in the literature: the construction
aimed at the viscous case does not use a layer cascade, but a self-similar core in which the
scale-local Reynolds number stays of order one.
\end{remark}

\section{A bound for the smooth-forcing design}\label{sec:design}

The inviscid Boussinesq and Euler constructions with smooth forcing fix their design as follows. The
frequencies satisfy $\lambda_{q}=\lambda_{q-1}^{Q_{q}}$ with $Q_{q}=Q^{\ast}+q$ and $Q^{\ast}\ge200$;
the number of derivatives of the force controlled uniformly is $k_{q}$, constrained by
$120k_{q}\le Q_{q}$; the localization error is cancelled through $J=2k+8$ correction levels, each
gaining a factor $\Pi^{5}\lambda^{-(1-\delta)}$ where $\Pi\le\lambda^{5/Q}$; and the seed is chosen so
that the remaining force stays below $\lambda^{-6}$. With the published amplitude rule a layer stops at
$\lambda^{-7/8}$, so the margin of (A1) is $\delta=1/8$.

\begin{theorem}\label{thm:design}
Under that design, carrying dissipation $-\nu(-\Delta)^{\alpha}$ requires
\[
J\Bigl[(1-\delta)-\frac5Q\Bigr]\ \ge\ k+6
\qquad\text{and}\qquad
\alpha\le\frac{\delta}{4Q},
\]
the first from the force budget and the second from~\eqref{eq:growth}. With $Q=120k$, the smallest
ratio the design permits, these give
\[
\delta\le1-\frac{k+6}{2k+8}-\frac1{24k},
\qquad
\alpha\le\frac{1}{480k}\Bigl(1-\frac{k+6}{2k+8}-\frac1{24k}\Bigr),
\]
whose maximum over $k\ge1$ is attained at $k=1$:
\[
\boxed{\ \alpha\le5.38\times10^{-4}\ \ \bigl(1.08\times10^{-3}\ \text{in the }|\nabla|^{\alpha}\ \text{convention}\bigr),\quad \delta\le0.2583,\ Q=120 .\ }
\]
\end{theorem}

\begin{proof}
Each correction level gains the factor $\Pi^{5}\lambda^{-(1-\delta)}$ and there are $J$ of them, so
the residual force carries the exponent $-J[(1-\delta)-5/Q]$ in $\lambda$ after using
$\Pi\le\lambda^{5/Q}$. Controlling $k$ derivatives of a field oscillating at frequency $\lambda$ costs
$\lambda^{k}$, and the design requires the result to stay below $\lambda^{-6}$, which is the first
display. The second is~\eqref{eq:growth} with $A_{q}=\lambda_{q-1}^{\delta}$ and
$\lambda_{q}=\lambda_{q-1}^{Q}$, taking the insertion angle of order one, which is the most favourable
case. Eliminating $\delta$ and setting $Q=120k$ gives the displayed formula; it is decreasing in $k$,
since the bracket increases to $1/2$ while the prefactor decreases like $1/k$, so the maximum is at
$k=1$.
\end{proof}

\begin{center}
\begin{tabular}{rrrrr}
\toprule
$k$ & $J=2k+8$ & $Q=120k$ & $\delta\le$ & $\alpha\le$ \\
\midrule
$1$ & $10$ & $120$ & $0.2583$ & $5.38\times10^{-4}$ \\
$2$ & $12$ & $240$ & $0.3125$ & $3.26\times10^{-4}$ \\
$3$ & $14$ & $360$ & $0.3433$ & $2.38\times10^{-4}$ \\
$10$ & $28$ & $1200$ & $0.4244$ & $8.84\times10^{-5}$ \\
$1000$ & $2008$ & $120000$ & $0.4990$ & $1.04\times10^{-6}$ \\
\bottomrule
\end{tabular}
\end{center}

The bound is $86$ times below the exponent $0.04633$ already reached with a force in
$L^{1}_{t}C^{1,\epsilon}_{x}\cap L^{\infty}_{t}L^{2}_{x}$, and the shortfall does not close by
retuning: at that exponent the ratio rule cannot buy even one derivative. Smoothness of the force and
fractional dissipation are purchased from a single budget in this design, with the frequency ratio as
the currency: $120$ units of ratio per derivative, and the exponent inversely proportional to the
ratio.

\begin{remark}\label{rem:levers}
Three quantities enter the proof, so a dissipative version of this design must move one of them:
decouple the correction levels from the frequency ratio, so that $Q$ need not grow with $k$; raise the
amplitude margin beyond $1/2$, which the hierarchy as stated forbids; or change the growth law, since
the factor four in $\alpha\le\delta/(4Q)$ comes from the rate $\sqrt A$ and becomes a factor two for
vortex stretching.
\end{remark}

\section{The exponent budget behind the published threshold}\label{sec:budget}

The hypodissipative Navier-Stokes construction writes every quantity as a power of one base: the
frequency of layer $n$ is $M_{n}=N^{R^{n}}$, the amplitude $A_{n}=N^{aR^{n}}$, the localization scale
$L_{n}=N^{bR^{n}}$, and $s>0$ is the H\"older margin of the force. Its force estimates give four
constraints:

\begin{center}
\begin{tabular}{llc}
\toprule
 & origin & constraint \\
\midrule
\textbf{D} & dissipation & $a=\alpha R$ \\
\textbf{S} & self-interaction & $2a-a/R+b+s-1\le0$ \\
\textbf{L} & localization & $b\ge a+1/R+s$ \\
\textbf{O} & outer velocity on the inner layer & $a+2/R+s+1-3b\le0$ \\
\bottomrule
\end{tabular}
\end{center}

\begin{proposition}\label{prop:threshold}
With \textbf{D} and \textbf{S} saturated, \textbf{O} reduces to $4s<2+3\alpha-7\alpha R-2/R$, whose
maximum over the frequency ratio is at $\alpha R^{2}=2/7$. At that optimum
$s(\alpha)=\bigl(2+3\alpha-2\sqrt{14\alpha}\bigr)/4$, and $s>0$ is equivalent to
$9\alpha^{2}-44\alpha+4>0$, whose smaller root is $(22-8\sqrt7)/9$.
\end{proposition}

The optimizing ratio $R=\sqrt{2/(7\alpha)}$ is the one the construction uses, and no coefficient in
the chain is fitted. At the threshold the localization constraint is inactive, with margin
$\tfrac12(\sqrt{2\alpha/7}-\alpha)=0.0350$; a budget that saturates \textbf{L} instead and omits
\textbf{O} yields $\alpha R^{2}=1/3$ and the weaker bound $5-2\sqrt6=0.10102$, the value given by an
order-of-magnitude estimate in the source. The difference between the two is precisely which
constraint is active.

\begin{proposition}\label{prop:ratios}
A positive force margin requires $7\alpha R^{2}-(2+3\alpha)R+2<0$, so the admissible ratios form the
interval with endpoints
$R_{\pm}=\bigl[(2+3\alpha)\pm\sqrt{9\alpha^{2}-44\alpha+4}\bigr]/(14\alpha)$. The discriminant is the
polynomial of Proposition~\ref{prop:threshold}, so the interval closes to a single point exactly at the
threshold, and that point is $\sqrt{2/(7\alpha)}$. Moreover $R_{-}>1$ for every $\alpha>0$, and at
$R=1$ the margin equals $-\alpha$.
\end{proposition}

\begin{center}
\begin{tabular}{lrr}
\toprule
$\alpha$ ($|\nabla|^{\alpha}$) & admissible $R$ & optimizing $R$ \\
\midrule
$0.001$ & $(1.002,\ 285.1)$ & $16.90$ \\
$0.01$ & $(1.021,\ 27.98)$ & $5.345$ \\
$0.05$ & $(1.143,\ 5.000)$ & $2.390$ \\
$0.09267$ & $\{1.7559\}$ & $1.7559$ \\
$0.10$ & empty & none \\
\bottomrule
\end{tabular}
\end{center}

\begin{corollary}
Cascades with geometrically growing frequencies, the case $R=1$, are inadmissible at every positive
dissipation: localization requires $b\ge a+1+s$ while self-interaction permits only $b\le1-a-s$,
forcing $\alpha+s\le0$. Super-geometric frequency growth is therefore necessary, and more of it is
necessary as the dissipation rises, until at the threshold a single ratio remains.
\end{corollary}

Propositions~\ref{prop:threshold} and~\ref{prop:ratios} are verified in exact arithmetic by computer
algebra in the accompanying repository.

\section{Verification of the dissipative extension}\label{sec:validation}

Theorems~\ref{thm:class} and~\ref{thm:design} rest on~\eqref{eq:diss}, which is derived rather than
taken from the constructions. It was tested against direct pseudospectral simulation of the nonlinear
two-dimensional Boussinesq system in double precision with two-thirds dealiasing.

\paragraph{The amplitude system.} On the stratification $\theta=-(A/\lambda_{0})\sin(\lambda_{0}x_{2})$
the measured peak growth rate matches $\sqrt A\sin\varphi$ to $4\times10^{-5}$ relative at two
insertion angles; the predicted independence of the rate from $\lambda$ holds to $2.7\times10^{-4}$
across an eightfold range; and the damping term of~\eqref{eq:diss} tracks the measured rate to
$2.6\times10^{-4}$ absolute across $\alpha\in\{1/4,1/2,1\}$ and $\nu\in\{10^{-5},10^{-4}\}$.

\paragraph{The cascade premise.} On a two-scale background that is an exact steady state, with drift
$1.45\times10^{-15}$, the local growth rate of a fine wave follows the total gradient at correlation
$0.99945$ and the base stratification alone at $0.212$: a layer responds to the accumulated
low-frequency gradient rather than to the base.

\paragraph{The growth, steering and holding cycle.} The control that returns a grown layer to rest
rotates the background gradient and the phase direction together, so that the laboratory component
$\zeta_{1}$ changes sign while $J\zeta\cdot G$ and $|\zeta|$ are preserved. In the co-rotating frame it
becomes a rotating gravity direction, with the stratification held by a low-frequency force. On a
background whose gradient is affine to fourth order the cycle reproduces the amplitude system to
$3.06\times10^{-6}$ in the steering gain and returns the vorticity amplitude to $3.67\times10^{-4}$ of
its peak, where it holds to $0.36\%$ over three growth times; on a plain sinusoidal background the
same quantities are $1.80\times10^{-2}$ and $1.57\times10^{-2}$, the difference being the curvature of
the background over the layer's support. Under dissipation the cycle factorizes as $e^{-dt}$ times the
inviscid cycle with $d=\nu\lambda^{2\alpha}$: a holding interval decays at $d$ to within $0.97\%$, and
the rescaled viscous run reproduces the inviscid one to $0.24\%$.

\paragraph{Localization.} Since the constructions localize each layer, \eqref{eq:diss} is applied to
a band rather than to a single frequency. For an envelope of radius $\ell$ on a wave of frequency
$\lambda$ the relative error of the damping law is
\[
\frac{\bigl\lVert(-\Delta)^{\alpha}(g\,e^{i\lambda\zeta\cdot x})-(\lambda|\zeta|)^{2\alpha}g\,e^{i\lambda\zeta\cdot x}\bigr\rVert}{\bigl\lVert(\lambda|\zeta|)^{2\alpha}g\bigr\rVert}
\;\approx\;\frac{3.6\,\alpha}{\ell\lambda},
\]
first order in the bandwidth ratio, with measured slope $-0.9863$, and depending on that ratio alone:
configurations sharing $\ell\lambda$ but splitting it differently between frequency and envelope agree
to four digits. The coefficient is linear in $\alpha$, as the symbol expansion predicts. At the
separations the constructions use, $\ell_{q}^{-1}=\lambda_{q-1}^{3}$ against
$\lambda_{q}=\lambda_{q-1}^{Q_{q}}$ with $Q_{q}\ge201$, the error is $\lambda_{q-1}^{3-Q_{q}}$, so
localization does not affect the bounds above.

\paragraph{A property of the mechanism in an unforced periodic setting.} The stratification these
layers grow on is Rayleigh-Taylor unstable at rate $\sqrt A$, while a layer inserted at angle $s$
grows at $\sqrt A\sin s$. Perturbations at more favourable angles therefore gain $1/\sin s$ times as
many e-foldings as the layer. The constructions are unaffected, since each layer is localized and the
force controls the exterior, but any realization without that control must keep its schedule short
enough for the layer to stay ahead.

\section{The status of the announced results}

The two September 2026 announcements are accompanied by machine-checked certificates. For the
Navier-Stokes and Euler development, the headline statements match the forced alternatives of the
official problem statement clause by clause; the development builds with no incomplete proofs and
its four headline theorems depend only on the three standard axioms; and every constant in the closure
of both root modules replays through the Lean kernel from an empty environment at the pinned toolchain
version. The reference statement is adapted from an independently maintained collection of formalized
conjectures, so the definitions that determine whether the formalized statement is the intended one
predate the claim. None of this speaks to the mathematics outside the formalization, and kernel replay
detects environment manipulation rather than serving as an external verifier.

\section{Scope}

The bounds of Sections~\ref{sec:class} and~\ref{sec:design} are necessary conditions for growth in
the layer mechanism, derived from the amplitude system, the amplification identity and the
dissipative term~\eqref{eq:diss}. They are exponent statements: constants, logarithmic factors and the
corrections that make the localized ansatz exact are not tracked, and the published inviscid results
are not affected by them, since those results carry no dissipation. The mechanism is due to C\'ordoba
and Mart\'inez-Zoroa, extended to smooth forcing by Alpog\'e and Buckmaster; nothing in this paper
establishes blowup for any equation or improves any published threshold.

An earlier version of this paper stated that the value $p=11/4+\sqrt7$, obtained by evaluating the
relation $\alpha_{c}=1/(4p)$ at the published threshold, carried information about that threshold. It
does not: the relation is constraint \textbf{D} saturated, which the construction saturates at every
$\alpha$, so the identity holds along the whole family, and its algebraic form is automatic because
$\alpha_{0}\in\Q(\sqrt7)$ and the norm of $22-8\sqrt7$ is a perfect square. That statement is
withdrawn.

\section*{Data and code}

The numerical data, the simulations, the exact-arithmetic verifications of
Propositions~\ref{prop:threshold} and~\ref{prop:ratios}, and the computations behind
Theorems~\ref{thm:class} and~\ref{thm:design} are public at \href{\recordurl}{the project repository}.

\begin{thebibliography}{9}
\bibitem{fefferman} C. L. Fefferman, \emph{Existence and smoothness of the Navier-Stokes equation},
Clay Mathematics Institute Millennium Prize problem description.
\bibitem{cmz} D. C\'ordoba, L. Mart\'inez-Zoroa and F. Zheng, \emph{Finite time blow-up for the
hypodissipative Navier Stokes equations with a force in $L^{1}_{t}C^{1,\epsilon}_{x}$ and
$L^{\infty}_{t}L^{2}_{x}$}, Arch. Ration. Mech. Anal. \textbf{250} (2026), article 38;
\href{https://arxiv.org/abs/2407.06776}{arXiv:2407.06776}.
\bibitem{clsmz} D. C\'ordoba, C. La\'in-Sanclemente and L. Mart\'inez-Zoroa, \emph{Blowup for the
Boussinesq equations with an infinite chain of degenerate pendula}, Adv. Math. \textbf{480} (2026);
\href{https://arxiv.org/abs/2505.20988}{arXiv:2505.20988}.
\bibitem{ab-boussinesq} L. Alpog\'e and T. Buckmaster, \emph{Blowup for the Boussinesq equations with
smooth forcing}, preprint, 2026.
\bibitem{ab-euler} L. Alpog\'e and T. Buckmaster, \emph{Blowup for the three-dimensional Euler
equations with smooth forcing}, preprint, 2026.
\bibitem{ab-ipm} L. Alpog\'e and T. Buckmaster, \emph{Blowup for the incompressible porous media
equation with smooth forcing}, preprint, 2026.
\bibitem{openai-ns} OpenAI, \emph{Finite-time blowup for forced three-dimensional Navier-Stokes},
manuscript, 2026.
\bibitem{openai-euler} OpenAI, \emph{Finite-time blowup for three-dimensional Euler}, manuscript,
2026.
\end{thebibliography}

\end{document}
